% Irodov Part V §5.3 Diffraction of Light, originals 97--108.
% Directly checked uploaded source pp.212--215 (printed 204--205) and answer p.329.
\diffractionproblem{97}{204}
单色平面波垂直照射一圆孔，孔径恰开放屏后距孔 $b$ 处 P 点的前 $N$ 个菲涅耳半波带，波长为 $\lambda$。
若已知屏面相对于 P 的距离 $r$ 的衍射光强分布 $I(r)$，求圆孔前入射光强 $I_0$。
\exsolution
平面波前 $N$ 个半波带的孔半径由
$R_N^2\simeq Nb\lambda$ 给出，孔面积
$A=\pi Nb\lambda$。
假设屏足够大收集全部透过光且无吸收，
总功率守恒：
\[
I_0A=2\pi\int_0^\infty I(r)r\,\mathrm dr.
\]
消去面积，得到
\[
\boxed{I_0=\frac2{Nb\lambda}
\int_0^\infty I(r)r\,\mathrm dr}.
\]
\excheck 一维径向积分须乘圆周因子 $2\pi r$，并假定衍射图样关于 P 轴对称。

\diffractionproblem{98}{204}
单色点光源位于圆孔屏前
$a=100$cm，波长 $\lambda=0.50\,\mathrm{\mu m}$，
孔半径 $r=1.0$mm。
求屏后观察点 P 到孔屏的距离 $b$，使孔内包含
$k=3$ 个菲涅耳半波带。
\exsolution
点源—孔—观察点几何中，
第 $k$ 个波带边缘半径满足
\[
r_k^2=k\lambda\frac{ab}{a+b}.
\]
将 $r=r_k$，
$r^2(a+b)=k\lambda ab$，
解出
\[
\boxed{b=\frac{ar^2}{k\lambda a-r^2}}.
\]
代 $a=1.0$m、$r=10^{-3}$m、
$k=3$、$\lambda=0.50\times10^{-6}$m，
\[
\boxed{b=2.0\,\mathrm m}.
\]
\excheck 分母必须为正 $k\lambda a>r^2$，否则有限正距离无解。

\diffractionproblem{99}{204}
可变半径圆孔屏距离单色点源
$a=100$cm、观察屏
$b=125$cm。孔半径从
$r_1=1.00$mm 增至
$r_2=1.29$mm 时，观察点中心由某强度极大转至\emph{紧接着}的强度极大。
求波长。
\exsolution
圆孔中心菲涅耳衍射的相邻强度极大，开放波带数相差2：
$N_2-N_1=2$。
波带关系
$r_N^2=N\lambda ab/(a+b)$，
故
\[
r_2^2-r_1^2=2\lambda\frac{ab}{a+b},
\qquad
\boxed{\lambda=
\frac{(r_2^2-r_1^2)(a+b)}{2ab}}.
\]
数值
\[
\boxed{\lambda\simeq
\frac{(1.29^2-1.00^2)10^{-6}(2.25)}
{2(1.00)(1.25)}
\simeq0.60\,\mathrm{\mu m}}.
\]
\excheck 相邻中心极大对应开放波带数差2而非1；开放波带数差1通常在亮暗之间转换。

\diffractionproblem{100}{204}
强度为 $I_0$ 的单色平面波垂直落到圆孔不透明屏。
观察点 P 所对应的圆孔：(a) 恰是第一菲涅耳波带，或仅为第一带内面积的一半；(b) 原开放完整第一带，后来沿直径挡住一半。
分别求 P 处的近似光强。
\exsolution
对完整无限波面，菲涅耳相邻半波带的复振幅近似相消。
第一带的复振幅约为无限自由波面振幅的2倍，
所以全开第一带时
\[
\boxed{I_{\text{第一带}}\simeq4I_0}.
\]
第一带内\emph{面积一半}对应径向相位跨度从0到 $\pi/2$，
复振幅为第一整带的矢量弦的
$1/\sqrt2$ 倍，
\[
\boxed{I_{\text{第一带内半面积}}\simeq2I_0}.
\]
若沿直径挡住第一带的半圆而不改变径向相位范围，
其复振幅恰缩小一半，
\[
\boxed{I_{\text{半圆孔}}\simeq I_0}.
\]
\excheck “第一带的内半面积”与“沿直径开放半圆”面积虽然相同，复振幅相位分布不同，因此光强不相同。

\diffractionproblem{101}{204}
不透明圆盘恰遮住观察点 P 的第一菲涅耳波带。
求在下列改变后 P 的近似光强：
(a) 沿直径移去该圆盘的一半；
(b) 只移去第一带外半面积的一半（沿直径半圆部分）。
入射平面波强度为 $I_0$。
\exsolution
把原完全开放波面的复振幅记 $A_0$；
第一带复振幅约为 $2A_0$，
其余无限外带和的振幅约为 $-A_0$。
(a) 移去一半圆盘后，该第一带半圆新贡献为
$(1/2)(2A_0)=A_0$，加外带 $-A_0$ 抵消，所以
\[
\boxed{I_a\simeq0}.
\]
(b) 第一带外半面积在径向相位 $\pi/2$ 到 $\pi$，
其复振幅相位相对第一半面积发生旋转。
移去该部分的一半圆周后，
复振幅与外带矢量和的模约为
$|A_0|/\sqrt2$，故
\[
\boxed{I_b\simeq\frac12I_0}.
\]
\excheck 沿径向切分与沿直径切分是不同的光路相量问题；原书结果采用菲涅耳波带法的近似。

\diffractionproblem{102}{204--205}
按原书图5.20的8种不透明屏形状，求点 P 处的光强。
(a) 屏1—3的顶角与屏4的半平面边缘均穿过 P；
(b) 屏5—8的圆弧边缘与第一菲涅耳波带边界重合。
\exsolution
设从 P 看被不透明扇区遮挡的张角为 $\varphi$。
(a) 在无限波面中去掉角比例 $\varphi/(2\pi)$ 后，
复振幅剩余比例
$1-\varphi/(2\pi)$，
\[
\boxed{I_a(\varphi)\simeq
\left(1-\frac{\varphi}{2\pi}\right)^2I_0}.
\]
图示屏1—4依次对应
$\varphi=\pi/2,\pi,3\pi/2,\pi$，因此
\[
\boxed{I_1:I_2:I_3:I_4
\simeq\frac9{16}:\frac14:\frac1{16}:\frac14}.
\]
(b) 圆弧边缘沿第一带外缘时，第一带扇区的附加复振幅约为
$(\varphi/2\pi)A_0$，
因此
\[
\boxed{I_b(\varphi)\simeq
\left(1+\frac{\varphi}{2\pi}\right)^2I_0}.
\]
图5.20屏5—8依次对应上述角度，得到
\[
\boxed{I_5:I_6:I_7:I_8
\simeq\frac{25}{16}:\frac94:\frac{49}{16}:\frac94}.
\]
\excheck 本题属于依图5.20编号的几何衍射题，角度必须与原图遮光扇区位置相匹配；结果保留了原书半波带近似。

\diffractionproblem{103}{205}
波长 $\lambda=0.60\,\mathrm{\mu m}$ 的平面波垂直入射厚玻璃板。
玻璃板背面有圆形凹槽，凹槽对观察点 P 对应第一又半个菲涅耳带（图5.21）。
求凹槽深 $h$，使 P 处强度
(a) 最大；(b) 最小；(c) 等于自由入射强度。
\exsolution
凹槽把局部玻璃替换为空气，相应区域附加相位
$\delta=2\pi(n-1)h/\lambda$（相位整体符号不影响强度）。
由第一又半波带的复振幅矢量和，再与未开槽波面相加，
干涉极值条件分别是
\[
\delta_a=2\pi\left(k+\frac38\right),\quad
\delta_b=2\pi\left(k+\frac78\right),
\]
以及使 $I=I_0$ 的两组相位
$\delta_c=2\pi k$ 或
$2\pi(k+\tfrac34)$。
故
\[
\boxed{h_a=\frac{\lambda}{n-1}
\left(k+\frac38\right)},\quad
\boxed{h_b=\frac{\lambda}{n-1}
\left(k+\frac78\right)},
\]
\[
\boxed{h_c=\frac{\lambda}{n-1}k
\quad\text{或}\quad
\frac{\lambda}{n-1}\left(k+\frac34\right)}.
\]
普通玻璃 $n\simeq1.50$ 时共同厚度周期为
$\lambda/(n-1)=1.2\,\mathrm{\mu m}$。
\excheck 三问是三个不同的\emph{复振幅相位条件}；
光强与相位不是线性函数。

\diffractionproblem{104}{205}
单色平面波 $\lambda$、强度 $I_0$ 垂直入射玻璃板，
板后方设圆孔，恰开放观察点 P 的第一菲涅耳带；
孔中间又刻出面积恰占该带一半的圆形凹槽。
求使 P 处强度最大的凹槽深度及最大强度。
\exsolution
第一菲涅耳带的内半面积、外半面积具有相等模的复振幅，
二者相位中心相差 $\pi/2$。
凹槽额外相位
$\delta=2\pi(n-1)h/\lambda$，
要使两部分复振幅完全同相，
应选择 $\delta$ 相对原初相差补偿为
$3\pi/2$（模 $2\pi$）：
\[
\boxed{h=\frac{\lambda}{n-1}
\left(k+\frac34\right)},\quad k=0,1,\ldots.
\]
第一带原合成模约为 $2|A_0|$，
两个半带各模约为 $\sqrt2|A_0|$；
最佳同相时总模
$2\sqrt2|A_0|$，因此
\[
\boxed{I_{\max}\simeq8I_0}.
\]
\excheck 相位补偿后各半带成为同相，而不是简单地将原第一带光强乘2。

\diffractionproblem{105}{205}
波长 $\lambda=0.57\,\mathrm{\mu m}$ 的平面光波照在折射率
$n=1.60$ 的玻璃圆盘上，
圆盘从观察点 P 看恰遮住第一又半个菲涅耳带。
求 P 点强度最大时玻璃盘的最小厚度，\emph{必须计入玻璃圆盘内部多次反射光的干涉}。
\exsolution
先将不放玻璃盘时 P 点的自由波复振幅记为 $A_0$。
在近似圆形的菲涅耳振动曲线上，第一又半个波带的场为
$F=(1+i)A_0$，其余未遮挡波带合成为 $O=-iA_0$，
故 $F+O=A_0$。
若盘的\emph{复振幅透射系数}相对于同厚度空气为 $\tau(h)$，
总场为
\[
 A_P=F\tau(h)+O
 =A_0\bigl[(1+i)\tau(h)-i\bigr].
\]
由于玻璃圆盘两表面均有菲涅耳反射，不能无条件设
$\tau=e^{i\delta}$、$|\tau|=1$。
令 $k_0=2\pi/\lambda$、单界面强度反射率
$R_f=[(n-1)/(n+1)]^2$。
对平行、无吸收玻璃片把全部往返透射的振幅作几何级数相加，
其相对于同厚度空气的透射振幅为
\[
 \tau(h)=
 \frac{(1-R_f)e^{i(n-1)k_0h}}
 {1-R_f e^{2ink_0h}}.
\]
因此
\[
 |\tau|^2
 =\frac{(1-R_f)^2}
 {(1-R_f)^2+4R_f\sin^2(nk_0h)}
 \le1,
\]
等号成立当且仅当玻璃片处于透射共振：
\[
 \boxed{2nh=m\lambda},\qquad m=1,2,\ldots .
\]
由三角不等式，
$|A_P/A_0|\le\sqrt2|\tau|+1\le\sqrt2+1$。
要使 P 点光强达到全局最大，须同时满足
\emph{盘的振幅全透射}以及\emph{盘内外两部分复振幅同相}。
处于上述透射共振时，$\tau=e^{i\delta}$，
$\delta=(n-1)k_0h$。
由于 $F$ 的相位是 $\pi/4$、$O$ 的相位是 $-\pi/2$，
同相要求
\[
 \frac\pi4+\delta\equiv-\frac\pi2\pmod{2\pi}
 \quad\Longleftrightarrow\quad
 \boxed{\delta=2\pi\left(k+\frac58\right)},
 \qquad k=0,1,\ldots .
\]
同时用 $2nh=m\lambda$ 消去厚度：
\[
 m=\frac{2n}{n-1}\left(k+\frac58\right)
 =\frac{16k+10}{3}.
\]
$m$ 必须为整数，$k=0,1$ 分别对应 $m=10/3,26/3$，
都不满足玻璃片自身的完全透射条件；
最小可行整数阶数是 $k=2$，此时 $m=14$。
于是
\[
 \boxed{h_{\min}
 =\frac{14\lambda}{2n}
 =\frac{\lambda}{n-1}\left(2+\frac58\right)
 =2.49375\,\mathrm{\mu m}\simeq2.5\,\mathrm{\mu m}}.
\]
此时 $\tau=e^{i5\pi/4}$，
$F\tau$ 与 $O$ 同相，故
\[
 \boxed{I_{\max}=(1+\sqrt2)^2I_0}.
\]
\excheck 若错误地忽略玻璃盘内部的多束透射干涉，
仅令单程光程差首次满足 $\delta=5\pi/4$，
将得到 $h=0.59375\,\mathrm{\mu m}$，但此时
$2nh/\lambda=10/3$，$|\tau|<1$，
无法达到全局最大强度。
本题原文特意要求计入盘内透射干涉；所以 $2.5\,\mathrm{\mu m}$
不是额外虚构的厚度下限，而是两项相位与透射条件的\emph{首次共同解}。

\diffractionproblem{106}{206}
波长 $\lambda=0.54\,\mathrm{\mu m}$ 的平面波穿过薄会聚透镜
$f=50$cm，透镜后紧靠可变半径圆孔；
距圆孔 $b=75$cm 处放观察屏。
求屏中心照度取极大的孔径半径。
\exsolution
透镜在孔平面产生二次波前相位
$-kr^2/(2f)$，
孔至屏传播又加
$kr^2/(2b)$。
等效二次相位系数
$k(1/b-1/f)r^2/2$；
相邻半波带边缘满足
$r_j^2|\frac1b-\frac1f|=j\lambda$。
对中心亮度，开放奇数个半波带时近似为极大，
故
\[
\boxed{r_j=
\sqrt{\frac{j\lambda fb}{|b-f|}}},
\qquad j=1,3,5,\ldots .
\]
代数值
\[
\boxed{r_j\simeq0.90\sqrt j\,\mathrm{mm}},
\quad j\ {\rm 为奇数}.
\]
\excheck 由于观察屏在焦平面后方，需用
$|1/b-1/f|$，不能使用自由空间传播半波带半径 $\sqrt{j\lambda b}$。

\diffractionproblem{107}{206}
平面单色光垂直入射圆孔，屏在孔后
$b=9.0$m处。
把孔直径缩小 $\eta=3.0$ 倍，
求新屏距 $b'$，使衍射图形与原图几何相似而所有线性尺寸都缩小 $\eta$ 倍。
\exsolution
在菲涅耳衍射中无量纲坐标组合为
$r^2/(\lambda b)$ 与
$R^2/(\lambda b)$，
其中 $R$ 为孔半径。
若孔半径和图样线性坐标均缩小为原来的 $1/\eta$，
要保持菲涅耳数不变须
$b'=b/\eta^2$。
故
\[
\boxed{b'=\frac b{\eta^2}
=\frac{9.0}{3^2}=1.0\,\mathrm m}.
\]
\excheck 不能将新屏距误写成 $b/\eta$；
菲涅耳衍射传播相位含的是横向尺度的\emph{平方}。

\diffractionproblem{108}{206}
不透明球直径 $D=40$mm，
位于波长 $\lambda=0.55\,\mathrm{\mu m}$ 的光源与照相底片间，
源到球距 $a=12$m，球到底片
$b=18$m。
求：(a) 横向尺寸为 $y=6.0$mm 的光源所形成的像宽；
(b) 球面随机不规则高度达到何量级时球真正有效遮光？
\exsolution
(a) 点源经不透明球几何影的横向投影，
扩展光源导致阴影半影放大：
\[
\boxed{y'=\frac ba y
=\frac{18}{12}(6.0\,\mathrm{mm})
=9.0\,\mathrm{mm}}.
\]
(b) 球边缘菲涅耳带的有效宽度量级由曲面边界的衍射展宽控制：
\[
\boxed{h_{\min}\sim\frac{ab\lambda}{D(a+b)}}.
\]
代数据
\[
\boxed{h_{\min}\simeq
\frac{12(18)(0.55\times10^{-6})}{0.040(30)}
\simeq1.0\times10^{-4}\,\mathrm m
=0.10\,\mathrm{mm}}.
\]
\excheck 原题(b)问随机表面粗糙度与球边缘菲涅耳带宽的\emph{量级}比较，不能视为材料吸收系数的精确阈值。
